\documentclass[11pt,a4paper]{article}

\usepackage[margin=1in]{geometry}
\usepackage{authblk}
\usepackage{booktabs}
\usepackage{array}
\usepackage{hyperref}
\usepackage{graphicx}
\usepackage{amsmath}

\title{Conventional-SC-Dataset: an auditable, element-centric literature and physical-parameter platform for superconductivity research}

\author[1]{[Yuan Ma, Qiang Guo]}
\affil[1]{[Your Affiliation]}
\date{}

\begin{document}
\maketitle

\begin{abstract}
Superconductor research increasingly depends on structured access to scattered information about superconducting systems: DOI-linked source papers, system names, symmetry and structural descriptions, reported transition temperatures, pressure conditions, and electron--phonon parameters. Public materials databases and recent superconductivity datasets have accelerated structure-property analysis, benchmark construction, and high-throughput screening, yet practical access to system-specific superconducting records is still often maintained through static spreadsheets, local notes, or project-specific scripts. Here we present Conventional-SC-Dataset, an auditable web platform that organizes superconductors and superconducting systems around element combinations, DOI-linked paper metadata, multiple physical-parameter records per paper, structural information, contributor identity, controlled visualization, and selectable data sources. The platform provides a periodic-table entry point, three standardized superconductor-search modes, system-specific browsing through DOI links, DOI-based paper submission, spreadsheet-based batch import of superconducting records, JSON-format superconducting-data import and export, an AI-screened literature database mode, HTSC-2025 benchmark browsing, Alexandria electron--phonon coupling endpoints, and a Tc-prediction page connected to uploaded structural and PDOS information. In the current repository snapshot, the codebase contains 15,572 lines across backend, frontend, and maintenance scripts, 68 backend routes, and nine main page templates. The main metadata database contains 118 elements and currently carries no paper records, whereas an AI-screened metadata database contains 688 papers, 1,010 compounds, and 2,624 physical-parameter records, including 1,921 records with a Tc field and 2,068 records with a pressure field; the bundled HTSC-2025 JSON fixture contains 140 benchmark-style records. This work does not claim a new superconducting mechanism or a validated predictive model. Instead, it establishes a reproducible data-management layer for converting fragmented superconductor records and adjacent computed datasets into reviewed, element-centric, reusable data.
\end{abstract}

\noindent\textbf{Keywords:} superconductivity database; conventional superconductors; hydrides; literature curation; physical parameters; data platform; scientific software

\section{Introduction}

The search for superconductors is now inseparable from data infrastructure. Conventional phonon-mediated superconductors can be analyzed with relatively mature first-principles workflows, while high-pressure hydrides, ambient-pressure candidates, and benchmark datasets have expanded the chemical and structural space that researchers must compare. Recent examples illustrate this shift. High-throughput calculations have charted Bardeen--Cooper--Schrieffer superconductors among experimentally known compounds \cite{bercx2025prxenergy}. Hydride-focused machine-learning studies have screened millions of candidate structures or built interpretable descriptors for reduced-pressure superconductivity \cite{pires2026ambient,chen2026descriptors}. HTSC-2025 provides a benchmark dataset for AI-driven critical-temperature prediction in ambient-pressure high-temperature superconductors \cite{han2025htsc}. HPCSD, the High-Pressure Crystal Structure Database, further demonstrates the importance of traceable, pressure-resolved data infrastructure for extreme-condition materials research \cite{wang2026hpcsd}.

These resources address essential but different parts of the superconductivity data problem. Structure databases and high-throughput archives prioritize crystal structures, energies, phase stability, or computed electron--phonon quantities. Benchmark datasets prioritize comparable model evaluation. Literature-extraction efforts prioritize transforming text into material-property records. However, a researcher often faces a more operational question before model training or mechanistic interpretation begins: for a selected element system, which papers have been reported, what physical parameters were entered from each paper, who contributed the record, which records have been reviewed, and which records are allowed to enter public charts? This question is not only a data-science problem; it is also a collaborative curation and review problem.

The bottleneck is especially visible for superconducting systems because one paper may report multiple compounds, pressure ranges, structures, samples, or calculated quantities. Reducing a paper to a single formula--Tc pair discards important context. Conversely, storing all evidence in free-form notes makes later filtering and comparison difficult. A useful platform must therefore separate paper-level metadata from data-point-level physical parameters, preserve structural and source-traceable information, and expose quality-control state inside the workflow rather than presenting it as a scientific contribution in itself.

Conventional-SC-Dataset was developed to meet this practical need. The platform is organized around a simple path: select elements, choose a data source, find the corresponding superconducting systems, browse DOI-linked papers or computed candidates, submit records with physical parameters, curate records in a review interface, and expose selected data through charts and statistics. This design is inspired by the infrastructure logic of recent database papers such as HPCSD: the central contribution is not a single prediction, but a standardized and traceable repository that makes fragmented research evidence easier to compare and reuse. The scientific object differs, however. HPCSD resolves pressure-dependent crystal structures and DFT-recomputed energetics, whereas Conventional-SC-Dataset resolves DOI-linked superconductor records, element combinations, structural descriptors, physical-parameter entries, and adapters to adjacent superconductivity datasets.

This article reports the current implementation and boundary of Conventional-SC-Dataset. We describe the data model, user workflow, system implementation, local data status, and relationship to existing superconductivity and materials databases. We also state what the present platform does not yet support: it is not a full-text literature-mining engine, not a validated Tc-prediction model, not a general social community, and not yet a public benchmark with a frozen versioned release. The richer AI-screened database used for snapshot statistics is treated here as local evidence of schema capacity and workflow use, not as a finalized public release unless a release policy, synchronization procedure, and persistent archive are added.

\section{Related data resources}

Materials informatics has benefited from several mature database paradigms. The Materials Project \cite{jain2013materialsproject}, OQMD \cite{kirklin2015oqmd}, AFLOW \cite{curtarolo2015aflow}, and NOMAD \cite{draxl2018nomad} show how standardized computational data can support large-scale discovery, while ICSD \cite{belsky2002icsd} and COD \cite{grazulis2012cod} provide important crystallographic foundations. These resources establish the expectation that materials data should be searchable, reusable, and connected to well-defined metadata.

Superconductivity-specific resources have developed along complementary directions. SuperCon and its subsequent machine-readable or ontology-oriented derivatives are central sources for historical superconducting-property data \cite{nimsSupercon,foppiano2023supercon2,ishii2023ontology}. The 3DSC dataset connects superconducting transition temperatures to crystal structures, addressing a key limitation of formula-only datasets \cite{sommer2023threedsc}. HTSC-2025 focuses on recent ambient-pressure high-temperature superconductors and is explicitly framed as a benchmark for AI-based Tc prediction \cite{han2025htsc}. High-throughput computational studies further provide curated candidate lists and calculated electron--phonon properties \cite{bercx2025prxenergy,pires2026ambient}.

HPCSD is particularly relevant as a writing and design reference because it frames a database as infrastructure for a fragmented field \cite{wang2026hpcsd}. HPCSD integrates experimental and theoretical high-pressure structures into a pressure-resolved repository, reporting 77,346 structural entries spanning 89 elements in its initial release. Its argument is that high-pressure structural data had remained sparse and disjointed compared with ambient-pressure materials databases. Conventional-SC-Dataset adopts a similar infrastructure stance but targets a different gap: superconducting systems and physical-parameter records require element-centric browsing, DOI-based source linkage, structural description, and human quality control before they can reliably support downstream statistics or model-building.

\section{Design principles}

Conventional-SC-Dataset follows five design principles that translate the curation gap into system constraints. The platform must begin from element combinations because this is a natural query unit for superconductivity researchers; it must separate paper metadata from physical-parameter entries because one paper can report several systems or conditions; it must record quality-control state because submitted records are not automatically publication-ready; it must keep prediction results distinct from reviewed literature records because the present evidence supports infrastructure rather than model validation; and it must treat external or computed datasets as explicit data-source modes rather than blending them silently into the reviewed local corpus.

\subsection{Element combinations as the primary entry point}

Superconductivity researchers frequently reason from chemical composition. The platform therefore starts from a periodic table instead of a free-text search box. Selected elements are normalized, sorted, and mapped into compound systems. Three search modes are supported: \texttt{only}, \texttt{combination}, and \texttt{contains}. The \texttt{only} mode retrieves records whose element set exactly matches the selected set; \texttt{combination} retrieves existing subsystems of the selected set; and \texttt{contains} retrieves larger systems containing the selected elements. This distinction allows the same interface to support precise binary or ternary searches and broader element-family exploration.

\subsection{Explicit data-source modes}

The compound page now separates four browsing modes: the reviewed local literature database, the AI-screened literature database, HTSC-2025, and Alexandria electron--phonon coupling data. The local and AI-screened modes return DOI-linked paper records with pagination and review-aware filters. The HTSC-2025 mode returns benchmark-style ambient-pressure superconducting candidates from a JSON fixture. The Alexandria mode is implemented as an adapter around a locally imported electron--phonon coupling database, exposing element search, minimum-Tc filtering, stability filtering, material detail, download, and CIF-style structure export endpoints. This source separation keeps provenance visible: literature records, benchmark candidates, and computed EPC entries are searchable from the same element workflow but are not treated as the same evidence type.

\subsection{Paper-level metadata separated from physical-parameter records}

The platform separates a literature record from its physical data points. The \texttt{papers} table stores DOI, title, authors, journal, year, abstract, contributor fields, quality-control status, review comments, and visualization eligibility. The \texttt{paper\_data} table stores one or more entries linked to a paper and a compound, including article type, superconductor type, chemical formula, crystal structure, Tc, pressure, electron--phonon coupling parameter, logarithmic phonon frequency, density of states at the Fermi level, sample name, data-source note, and sequence in the paper. This separation preserves the fact that a single publication can contain multiple superconducting systems or multiple pressure--Tc records.

\subsection{Curation as a reviewed workflow}

Data entry is treated as a workflow rather than a one-step upload. Ordinary users can register, verify email, log in, and submit papers from a compound page. Curators can review and edit submissions, and privileged maintainers can manage user permissions. Review status, reviewer identity, review time, review comment, and visualization eligibility are stored with the paper record. These fields are quality-control infrastructure; they are not presented as standalone scientific innovations.

\subsection{Prediction as a structure-input tool}

The repository contains a Tc-prediction page for structure and PDOS inputs. This module is designed as a structure-input tool that reads uploaded CONTCAR and PDOS files, including hydrogen-projected PDOS information, and returns an estimated Tc-related result for exploratory use. It does not replace literature curation and is not presented here as a validated predictive contribution. This boundary is important because current evidence supports the platform as data infrastructure, not as a benchmarked machine-learning model.

\section{System implementation}

The implementation combines frontend pages, backend APIs, relational data objects, data-source adapters, and maintenance scripts. This section therefore describes the system at the level needed to reproduce the curation workflow rather than at the level of individual user-interface controls.

\subsection{Data model}

The current data model contains six central object types: elements, compounds, users, papers, physical-parameter records, and visual attachments. Elements store the 118 chemical elements and support the periodic-table interface. Compounds store chemical formulas and normalized element lists. Users store ordinary-user, curator, and maintainer states. Papers store literature metadata and quality-control fields. Physical-parameter records store per-paper compound and superconductivity data. Visual-attachment records currently support source images and can be evolved toward generated structure visualizations.

The implementation also contains non-curation data adapters. The AI-screened paper API reuses the paper and physical-parameter response model against a separate metadata database. HTSC-2025 is loaded from a repository JSON file and exposed through search, statistics, and detail endpoints. Alexandria support expects a locally built SQLite database under the Alexandria data directory and stores heavy raw fields outside normal literature tables; API responses deliberately return lightweight material summaries, structure metadata, and generated CIF text instead of large force-constant or dynamical-matrix payloads.

\begin{table}[t]
\centering
\caption{Core database objects and their manuscript-level roles.}
\label{tab:model}
\begin{tabular}{>{\raggedright\arraybackslash}p{0.18\linewidth} >{\raggedright\arraybackslash}p{0.34\linewidth} >{\raggedright\arraybackslash}p{0.36\linewidth}}
\toprule
Object & Main fields & Role in the platform \\
\midrule
\texttt{elements} & Symbol, English name, Chinese name, atomic number & Periodic-table rendering and element validation \\
\texttt{compounds} & Chemical formula, element list, composition, element IDs, element ratio & Element-combination search and compound grouping \\
\texttt{papers} & DOI, title, authors, journal, year, abstract, contributor fields, review fields, chart flag & Literature-level metadata and review state \\
\texttt{paper\_data} & Article type, superconductor type, formula, structure, Tc, pressure, $\lambda$, $\omega_{\log}$, $N(E_F)$, sample note & Multiple physical-parameter records per paper \\
\texttt{paper\_images} & Paper ID, figure labels, original images, thumbnails, file size & Visual attachments and source/structure context \\
\texttt{users} & Email, real name, password hash, email verification, admin flags, approval state & Authentication, submission, curation, and permission control \\
\bottomrule
\end{tabular}
\end{table}

\subsection{User-facing workflow}

The main user workflow begins at the home page. A user selects elements from the periodic table, chooses a search mode, and enters a compound page. The compound page loads superconducting-system information and retrieves records through backend APIs according to the selected data source. Users can filter local and AI-screened literature records by quality-control status, keyword, year range, and pagination; Alexandria adds minimum-Tc and stability filters; HTSC-2025 exposes benchmark-candidate pagination and detail views. Each local or AI-screened paper card can display metadata, DOI, authors, article type, superconductor type, chemical formula, crystal structure, abstract, and physical-parameter entries.

Authenticated users can submit a paper from a compound page. The submission requires DOI information, element context, article type, superconductor type, and at least one physical-parameter entry. The current implementation can store visual attachments, while the intended display direction is to emphasize crystal-structure visualization rather than screenshots as a manuscript-level contribution. The backend validates login state, parses JSON-formatted element and physical-parameter payloads, checks DOI metadata, creates or retrieves a compound, prevents duplicate DOI entries within a compound system, creates the paper, and writes physical-parameter records.

\subsection{Administrator workflow}

The curator interface supports the quality-control side of the platform. Curators can inspect unreviewed papers, update review status, edit records, and decide whether a paper should appear in homepage charts. Maintainers can approve curator applications and update user permissions. This role hierarchy is supporting infrastructure for trustworthy data management rather than a headline contribution of the manuscript.

\subsection{Import, export, and maintenance}

The platform also supports batch-oriented data operations. Here, batch import means importing multiple superconducting paper and physical-parameter records at once from an agreed spreadsheet or text format, rather than entering each DOI record manually. Maintenance scripts support database initialization, JSON-format superconducting-data export and import, ID migration, super-administrator creation, and backup-oriented operations. These scripts are part of the platform's reproducibility story: the database is not only a website state, but a data object that can be initialized, exported, imported, and maintained.

\section{Current repository and data status}

We evaluated the repository snapshot after synchronizing the \texttt{mayuan} branch with \texttt{origin/master} on 8 June 2026. The codebase contains 15,572 lines across Python, JavaScript, HTML, CSS, and shell scripts when backend, frontend, and maintenance scripts are counted. It includes 68 backend routes and nine main page templates. The current \texttt{backend/api} directory contains nine non-\texttt{\_\_init\_\_} API modules. These numbers are implementation descriptors rather than scientific performance metrics, but they indicate that the platform is beyond a static spreadsheet or isolated prototype script.

For journal submission, these snapshot descriptors should be regenerated from an archived revision and reported together with the exact counting protocol. The reproducibility package should include: the repository commit or release tag; the file extensions included in line-count statistics; the route-count command over FastAPI decorators; the database files used for Table~\ref{tab:status}; and a fixture-based workflow check that initializes a database, imports a small paper set, exports it, re-imports it, and compares record counts and key fields. The same package should document required submission fields, DOI metadata failure behavior, duplicate-DOI rules, review-state transitions, visualization eligibility, and the planned structure-visualization path. These steps are not model-performance experiments; they are the minimum reproducibility protocol for a database-system contribution.

\begin{table}[t]
\centering
\caption{Implementation and local data status in the current repository snapshot.}
\label{tab:status}
\begin{tabular}{>{\raggedright\arraybackslash}p{0.42\linewidth} r >{\raggedright\arraybackslash}p{0.34\linewidth}}
\toprule
Metric & Value & Scope \\
\midrule
Total code lines & 15,572 & Backend, frontend, and maintenance scripts \\
Backend routes & 68 & FastAPI route decorators in the backend \\
Main page templates & 9 & HTML templates under the frontend templates directory \\
Elements in \texttt{hydride\_literature.db} & 118 & Main metadata database \\
Papers in \texttt{hydride\_literature.db} & 0 & Main metadata database, current snapshot \\
Compounds in \texttt{hydride\_literature\_ai.db} & 1,010 & AI-screened metadata database \\
Papers in \texttt{hydride\_literature\_ai.db} & 688 & AI-screened metadata database \\
Distinct DOI values in \texttt{hydride\_literature\_ai.db} & 429 & Non-empty DOI entries \\
Physical-parameter records in \texttt{hydride\_literature\_ai.db} & 2,624 & \texttt{paper\_data} table \\
Records with Tc field & 1,921 & Non-empty \texttt{paper\_data.tc} \\
Records with pressure field & 2,068 & Non-empty \texttt{paper\_data.tc\_press} \\
Year range in AI-screened papers & 1909--2025 & Minimum and maximum paper year \\
Journal count in AI-screened papers & 105 & Distinct non-empty journal names \\
HTSC-2025 fixture records & 140 & Bundled JSON benchmark-style candidate records \\
\bottomrule
\end{tabular}
\end{table}

The current data status should be interpreted carefully. The main metadata database, \texttt{hydride\_literature.db}, contains the element dictionary and one compound record but no paper or physical-parameter records in this snapshot. A separate AI-screened metadata database, \texttt{hydride\_literature\_ai.db}, contains substantially more literature and physical-parameter data. The HTSC-2025 fixture is bundled as a JSON data source rather than as reviewed local literature records. Alexandria endpoints are implemented, but they require a locally imported Alexandria SQLite database to be present before meaningful row counts can be reported. Therefore, the platform's data model and workflows can support a non-trivial literature corpus and external computed-data browsing, but the manuscript should distinguish between the default main database, AI-screened data file, bundled benchmark fixture, and optional imported EPC database until a formal public release or synchronization policy is established.

\section{Comparison with recent datasets}

The closest conceptual comparison is not a single dataset but a family of infrastructure efforts. HPCSD standardizes high-pressure crystal structures and pressure-resolved energetics, with two data streams from reported elemental phases and CALYPSO crystal-structure prediction tasks \cite{wang2026hpcsd}. HTSC-2025 curates a benchmark of ambient-pressure high-temperature superconductors for AI-driven Tc prediction \cite{han2025htsc}. The PRX Energy BCS-superconductor study uses high-throughput first-principles calculations to identify promising phonon-mediated superconductors among experimentally known compounds \cite{bercx2025prxenergy}. Hydride machine-learning studies expand candidate discovery through large-scale screening or interpretable descriptors \cite{pires2026ambient,chen2026descriptors}.

\begin{table}[t]
\centering
\caption{Functional positioning of Conventional-SC-Dataset relative to adjacent resources.}
\label{tab:comparison}
\begin{tabular}{>{\raggedright\arraybackslash}p{0.18\linewidth} >{\raggedright\arraybackslash}p{0.22\linewidth} >{\raggedright\arraybackslash}p{0.22\linewidth} >{\raggedright\arraybackslash}p{0.25\linewidth}}
\toprule
Resource family & Primary object & Main strength & Complementary gap addressed here \\
\midrule
Materials databases & Computed materials entries, structures, energies, metadata & Large-scale reusable materials data and interoperability & Superconductor-specific DOI linkage, system records, physical-parameter rows, and visualization eligibility \\
Structure databases & Crystallographic structures & Curated or open crystal-structure reference data & Literature-level superconductivity parameters and structure-aware browsing workflow \\
SuperCon-style resources & Superconducting formula and property records & Historical superconductivity property coverage & Superconductor-search modes, per-paper multiple records, DOI links, and web curation flow \\
3DSC and HTSC-2025 & Structure-augmented or benchmark superconductivity datasets & Structure connection and AI-facing benchmark scope & Live curation workflow before frozen benchmark construction \\
HPCSD & Pressure-resolved high-pressure structures & Traceable structure repository for extreme-condition materials & Literature and physical-parameter curation rather than DFT-reoptimized structure infrastructure \\
\bottomrule
\end{tabular}
\end{table}

Conventional-SC-Dataset is deliberately positioned below these tasks in the data pipeline. It does not recompute all structures with a unified DFT protocol as HPCSD does. It does not define a frozen benchmark split as HTSC-2025 does, although it can browse a bundled HTSC-2025-style fixture. It does not claim exhaustive high-throughput electron--phonon calculations, although it now provides an Alexandria adapter for imported EPC records. Instead, it supports the curation layer that many downstream tasks require: mapping superconducting systems to DOI-linked papers, recording multiple physical parameters per paper, preserving structural descriptors, and controlling which records are suitable for visualization and reuse. In this sense, the platform is complementary to structure databases, benchmark datasets, and high-throughput computational screens.

\section{Limitations and next steps}

The present platform has several limitations. Data-model limitations include the fact that the relationship between a submitted paper and a registered user is currently represented mainly by contributor name and affiliation fields, not by a stable contributor-user foreign key. Storage and maintenance limitations include BLOB-based visual-attachment storage, the planned transition toward generated crystal-structure visualization, and the absence of a formal schema-migration system such as Alembic. Workflow and search limitations include the fact that pressure and Tc fields are available in physical-parameter records but are not yet elevated into fully mature first-class filtering dimensions across every local search path, even though Alexandria and HTSC-oriented modes already expose source-specific filters. Release and data-source limitations include the lack of a frozen public dataset release, persistent archive identifier, documented synchronization policy between the main and AI-screened databases, and a bundled imported Alexandria database with reproducible import metadata. Community limitations include the absence of likes, comments, bullet-screen messages, heat ranking, and database snapshot management.

These limitations define a practical development path. The next version should prioritize a stable \texttt{contributor\_user\_id} relationship, pressure and Tc range filters, clearer indexing strategy, backend-managed citation export, and formal schema migration. After these foundations are stable, the platform can more safely add public release versioning, dataset cards, benchmark splits, or community-interaction features.

\section{Conclusion}

Conventional-SC-Dataset provides an auditable, element-centric platform for superconducting-system and physical-parameter curation. By combining periodic-table entry, standardized superconductor-search modes, DOI-based paper linkage, multiple physical-parameter records per paper, structural descriptors, quality-control state, visualization eligibility, JSON-format superconducting-data import/export utilities, AI-screened browsing, HTSC-2025 fixture browsing, and Alexandria EPC adapters, it converts fragmented records and adjacent computed datasets into a maintainable data workflow. Its current contribution is infrastructural: it supports reviewed data organization for superconductivity research and prepares the ground for more reliable statistics, model-training corpora, and benchmark construction. Future work should focus on formalizing user--paper relationships, improving physical-parameter filtering, versioning public data releases, and integrating migration and synchronization policies before stronger claims about prediction or community-scale deployment are made.

\section*{Data and code availability}

The source code and local database files are available in the project repository at \texttt{[repository URL or local project path]}. A formal public dataset release, version identifier, and persistent archive DOI should be added before journal submission. The current manuscript statistics were computed after synchronizing the \texttt{mayuan} branch with \texttt{origin/master} on 8 June 2026.

\section*{Author contributions}

\textbf{[Your Name]} designed and implemented the platform, curated the manuscript scope, and prepared the article draft. Additional contributor roles should be filled in after confirming collaborators, code contributors, and data curators.

\section*{Competing interests}

The author declares no competing interests. \textbf{[Revise if needed before submission.]}

\section*{Acknowledgements}

The author thanks the developers and authors of the public superconductivity, materials, and high-pressure database studies cited in this manuscript. \textbf{[Add grant numbers, collaborators, and institutional support before submission.]}

\begin{thebibliography}{99}

\bibitem{jain2013materialsproject}
A. Jain, S. P. Ong, G. Hautier, W. Chen, W. D. Richards, S. Dacek, S. Cholia, D. Gunter, D. Skinner, G. Ceder, and K. A. Persson.
The Materials Project: A materials genome approach to accelerating materials innovation.
\textit{APL Materials} \textbf{1}, 011002 (2013).
doi:10.1063/1.4812323.

\bibitem{kirklin2015oqmd}
S. Kirklin, J. E. Saal, B. Meredig, A. Thompson, J. W. Doak, M. Aykol, S. Ruehl, and C. Wolverton.
The Open Quantum Materials Database (OQMD): assessing the accuracy of DFT formation energies.
\textit{npj Computational Materials} \textbf{1}, 15010 (2015).
doi:10.1038/npjcompumats.2015.10.

\bibitem{curtarolo2015aflow}
S. Curtarolo, W. Setyawan, G. L. W. Hart, M. Jahnatek, R. V. Chepulskii, R. H. Taylor, S. Wang, J. Xue, K. Yang, O. Levy, M. J. Mehl, H. T. Stokes, D. O. Demchenko, and D. Morgan.
AFLOW: The automatic flow framework for materials discovery.
\textit{Computational Materials Science} \textbf{58}, 218--226 (2012).
doi:10.1016/j.commatsci.2012.02.005.

\bibitem{draxl2018nomad}
C. Draxl and M. Scheffler.
NOMAD: The FAIR concept for big-data-driven materials science.
\textit{MRS Bulletin} \textbf{43}, 676--682 (2018).
doi:10.1557/mrs.2018.208.

\bibitem{belsky2002icsd}
A. Belsky, M. Hellenbrandt, V. L. Karen, and P. Luksch.
New developments in the Inorganic Crystal Structure Database (ICSD): accessibility in support of materials research and design.
\textit{Acta Crystallographica Section B} \textbf{58}, 364--369 (2002).
doi:10.1107/S0108768102006948.

\bibitem{grazulis2012cod}
S. Grazulis, A. Daskevic, A. Merkys, D. Chateigner, L. Lutterotti, M. Quiros, N. R. Serebryanaya, P. Moeck, R. T. Downs, and A. Le Bail.
Crystallography Open Database (COD): an open-access collection of crystal structures and platform for world-wide collaboration.
\textit{Nucleic Acids Research} \textbf{40}, D420--D427 (2012).
doi:10.1093/nar/gkr900.

\bibitem{nimsSupercon}
National Institute for Materials Science.
SuperCon Datasheet, version v240322.
doi:10.48505/nims.4487.

\bibitem{foppiano2023supercon2}
L. Foppiano, T. C. T. Ting, J. M. Romary, and coauthors.
Automatic extraction of materials and properties from superconductors scientific literature.
\textit{Science and Technology of Advanced Materials: Methods} \textbf{3}, 2153633 (2023).
doi:10.1080/27660400.2022.2153633.

\bibitem{ishii2023ontology}
Y. Ishii and T. Sakamoto.
Structuring superconductor data with ontology: reproducing historical datasets as knowledge bases.
\textit{Science and Technology of Advanced Materials: Methods} \textbf{3}, 2223051 (2023).
doi:10.1080/27660400.2023.2223051.

\bibitem{sommer2023threedsc}
T. Sommer, D. Di Cataldo, and L. Boeri.
3DSC -- a dataset of superconductors including crystal structures.
\textit{Scientific Data} \textbf{10}, 816 (2023).
doi:10.1038/s41597-023-02721-y.

\bibitem{han2025htsc}
X.-Q. Han, Z.-F. Gao, X.-D. Wang, Z. Ouyang, P.-J. Guo, and Z.-Y. Lu.
HTSC-2025: A benchmark dataset of ambient-pressure high-temperature superconductors for AI-driven critical temperature prediction.
\textit{Chinese Physics B} \textbf{34}, 100301 (2025).
doi:10.1088/1674-1056/adf042.

\bibitem{bercx2025prxenergy}
M. Bercx, S. Ponce, Y. Zhang, G. Trezza, A. G. Ghezeljehmeidan, L. Bastonero, J. Qiao, F. O. von Rohr, G. Pizzi, E. Chiavazzo, and N. Marzari.
Charting the landscape of Bardeen-Cooper-Schrieffer superconductors in experimentally known compounds.
\textit{PRX Energy} \textbf{4}, 033012 (2025).
doi:10.1103/sb28-fjc9.

\bibitem{pires2026ambient}
P. R. Pires, T. H. B. da Silva, K. Gao, K. H. Hiorth, T. F. T. Cerqueira, T. Cavignac, P.-P. De Breuck, H.-C. Wang, D. Dangic, Y.-W. Fang, A. Sanna, W. Cui, I. Errea, P. Torma, and M. A. L. Marques.
Machine learning driven exploration of hydride superconductors at ambient pressure.
\textit{Computational Materials Today} \textbf{10}, 100052 (2026).
doi:10.1016/j.commt.2026.100052.

\bibitem{chen2026descriptors}
J. Chen, J. Peng, Y. Liang, R. Wang, H. Dong, and W. Zhang.
Interpretable descriptors enable prediction of hydrogen-based superconductors at moderate pressures.
\textit{Materials Today Physics} \textbf{63}, 102073 (2026).
doi:10.1016/j.mtphys.2026.102073.

\bibitem{wang2026hpcsd}
Z. Wang, Q. Wang, J. Duan, H. Ge, X. Luo, P. Gao, W. Zhang, J. Lv, Y. Wang, and Y. Ma.
High-Pressure Crystal Structure Database.
\textit{arXiv} 2605.14471 (2026).

\end{thebibliography}

\end{document}
